% paper/sections/11b_conservation.tex
%
% §11.2  保存則の検証（Conservation Properties）
% ═══════════════════════════════════════════════════════════════════════════════
\subsection{保存則の検証}
\label{sec:ns_conservation}
% ═══════════════════════════════════════════════════════════════════════════════

NS ソルバが物理的に正しい保存則を再現することを，
Taylor--Green 渦の減衰問題により検証する．

% ─────────────────────────────────────────────────────────────────────────────
\subsubsection{運動エネルギー減衰}
\label{sec:energy_conservation}
% ─────────────────────────────────────────────────────────────────────────────

\textbf{設定：}
計算領域 $[0, 2\pi]^2$，周期境界条件，単相流（$\rho = 1$，$\nu = 0.01$，$\mathit{Re} = 100$）．
初期条件は Taylor--Green 渦：
\begin{align}
  u(x,y,0) &= \sin x\,\cos y, &
  v(x,y,0) &= -\cos x\,\sin y, \\
  p(x,y,0) &= -\tfrac{1}{4}(\cos 2x + \cos 2y).
\end{align}
厳密解は $\bu(t) = \bu(0)\,e^{-2\nu t}$ であり，
運動エネルギー $E_k(t) = E_k(0)\,e^{-4\nu t}$ が既知である．

$N = 64$，$\Delta t = 0.01$ で AB2+IPC スキームにより $T = 2.0$ まで計算し，
$E_k(t)$ の理論減衰率に対する相対誤差と
$\|\bnabla\cdot\bu\|_\infty$ の時間推移を記録する（表~\ref{tab:tgv_energy}）．

\begin{table}[htbp]
\centering
\caption{Taylor--Green 渦の運動エネルギー保存性と非圧縮条件（$N=64$，$\mathit{Re}=100$）}
\label{tab:tgv_energy}
\begin{tabular}{rcccc}
\toprule
$t$ & $E_k$（数値） & $E_k$（解析） & 相対誤差 & $\|\bnabla\cdot\bu\|_\infty$ \\
\midrule
$0.5$ & $9.674173 \times 10^{0}$ & $9.674173 \times 10^{0}$ & $7.1 \times 10^{-11}$ & $9.8 \times 10^{-14}$ \\
$1.0$ & $9.482612 \times 10^{0}$ & $9.482612 \times 10^{0}$ & $1.4 \times 10^{-10}$ & $9.7 \times 10^{-14}$ \\
$1.5$ & $9.294843 \times 10^{0}$ & $9.294843 \times 10^{0}$ & $2.1 \times 10^{-10}$ & $9.5 \times 10^{-14}$ \\
$2.0$ & $9.110793 \times 10^{0}$ & $9.110793 \times 10^{0}$ & $2.8 \times 10^{-10}$ & $9.2 \times 10^{-14}$ \\
\bottomrule
\end{tabular}
\end{table}

\textbf{結果：}
運動エネルギーの最大相対誤差は $2.8 \times 10^{-10}$（$t = 2.0$）であり，
目標の $10^{-4}$ を 6 桁以上下回る．
CCD-PPE による投影ステップは各タイムステップで
$\|\bnabla\cdot\bu\|_\infty < 10^{-13}$ を維持しており，
離散非圧縮条件が計算機精度水準で成立する．

% ─────────────────────────────────────────────────────────────────────────────
\subsubsection{離散非圧縮条件の定量評価}
\label{sec:div_free}
% ─────────────────────────────────────────────────────────────────────────────

上記 Taylor--Green 渦シミュレーション全 200 ステップを通じて，
$\|\bnabla\cdot\bu\|_\infty$ は $\Ord{10^{-13}}$--$\Ord{10^{-14}}$ に維持される．
これは CCD 微分演算子の 6 次精度と Spectral PPE ソルバの組み合わせにより，
速度場の無発散補正が浮動小数点丸め誤差のレベルで実現されていることを示す．
